\chapter{기계 디버깅}
\chaptermark{기계 디버깅}
\label{sec:debugging}

\section{회로도 없이 기계를 디버깅하는 법}

회로도 없이 작동하는 보드를 받은 엔지니어는 네 가지 기법으로 디버깅한다. \emph{프로브}를 대고 전선에 무엇이 흐르는지 본다. \emph{시험 신호를 넣고} 무엇이 바뀌는지 본다. 회로의 일부를 \emph{끊고} 무엇이 멈추는지 본다. 그리고 \emph{제어 레지스터를 읽어} 보드가 어떤 모드로 동작하는지 알아낸다. 이 책 전체가 뇌의 회로도를 읽으려는 시도다. 이 장에서는 네 기법 가운데 엔지니어가 이미 쓸 줄 아는 것이 무엇이고, 앞 장들이 무엇을 기대하라고 말해 주는지 살펴본다.

오늘날 이런 디버깅의 가장 눈에 띄는 예는 뇌-컴퓨터 인터페이스다. 뉴런의 스파이크를 읽어 외부 기계를 위한 명령으로 바꾸거나, 반대로 바깥에서 뇌로 신호를 넣는 장치다. 이 장의 대부분은 이 인터페이스, 특히 Neuralink가 하는 일을 다룬다.

\section{프로브: 전극은 무엇을 듣는가}

뇌의 프로브는 조직에 꽂은 가느다란 전극이다. 전극은 반경 약 100마이크로미터 안에 있는 뉴런, 보통 한 개에서 몇 개의 스파이크를 듣는다. 전극에 잡힌 스파이크는 수십에서 수백 마이크로볼트, 길이 약 1밀리초의 파형이며, 그 모양으로 이웃 뉴런을 서로 구별할 수 있다. 가장 잘 들리는 것은 큰 \nt{GLUT} 뉴런이다. 피질에서 대다수를 차지하고(표~\ref{tab:types}) 전류도 더 세다.

주된 어려움은 바로 보인다. 오늘날 좋은 프로브는 전극 수백에서 수천 개이지만, 뇌의 뉴런은 $8.6\cdot10^{10}$개다. 우리는 기계의 대략 1억분의 1을 엿듣는다. 그렇게 보잘것없는 표본에서 어떻게 무언가를 알 수 있을까? 답은 \ref{sec:code}장에 있다. 이웃한 뉴런들의 잡음은 서로 상관되어 있어서, 무리에 대한 평균은 한계에 부딪힌다(명제~\ref{prop:pooling}). 뉴런 100개는 1000개와 거의 같은 양을 전달한다. 게다가 운동 피질이 푸는 문제는 저차원이다. 팔의 관절은 많지 않고(\ref{sec:muscles}장), 팔을 제어하는 수백 개 뉴런의 활동은 공통 변수 열 개 남짓으로 이루어진 공간에 들어간다~\cite{Gallego2017}. 뉴런 100개짜리 프로브는 이 변수를 거의 다 듣는다. 만약 뇌가 뉴런마다 독립된 메시지를 저장한다면 이런 엿듣기는 쓸모가 없을 것이다.

\section{데이터 흐름: 프로브에서 바로 압축하기}

프로브는 엄청난 흐름을 낸다. 스파이크의 모양을 보려면 각 전극을 초당 약 $2$만 번, 약 10비트 정밀도로 디지털화해야 한다. 전극 1000개는
\begin{empheq}[box=\keybox]{equation}
\begin{aligned}
R_{\text{원시}} \;&=\; N\,f_s\,b \;\approx\; 10^3\cdot2\cdot10^4\cdot10 \;\approx\; 2\cdot10^8\ \text{bit/s},\\
R_{\text{스파이크}} \;&\approx\; N\,r\,\log_2\frac{e}{r\,\Delta t} \;\approx\; 8\cdot10^4\ \text{bit/s}
\end{aligned}
\label{eq:probedata}
\end{empheq}
를 낸다. 두 번째 추정은 $r=10$~spikes/s, 정밀도 $\Delta t=1\ms$에서의 스파이크 흐름의 용량~\eqref{eq:spikeentropy}이다. 흐름에 실제로 들어 있는 것은 2500분의 1 크기다. 이식 장치에는 이것이 결정적인 차이다. 머리 안에서 조직을 데우지 않고 쓸 수 있는 전력으로는 원시 흐름을 피부를 통해 무선으로 보낼 수 없다. 그래서 이식 장치는 전극 옆 칩에서 바로 스파이크를 골라내고, 이벤트, 즉 채널 번호와 스파이크 시각만 밖으로 보내야 한다. Neuralink 시스템을 처음 소개한 논문은 아직 거기까지 가지 못했다. 칩이 신호를 증폭하고 디지털화하면 원시 흐름 전체가 케이블로 나가고, 스파이크는 외부 프로그램이 골라낸다. 칩 위의 압축과 무선 통신은 계획으로 언급되어 있다~\cite{Musk2019}.

익숙한 기법이다. 눈은 긴 케이블로 보내기 전에 신호를 100분의 1로 압축한다(\ref{sec:sensors}장). 이벤트 카메라는 프레임이 아니라 변화를 보낸다. 프로브도 전선에 맞추려면 뇌 자신이 하는 일을 해야 한다. 스파이크로, 그리고 새 소식만 말하는 것이다.

\section{Neuralink가 하는 일}

Neuralink 장치는 이런 제약에 맞춰 만든 프로브다~\cite{Musk2019}. 딱딱한 바늘 배열 대신, 머리카락보다 가는 유연한 실 여러 가닥에 전극이 줄지어 달려 있다. 시스템을 처음 소개한 논문에서는 실 $96$가닥에 전극 최대 $3072$개, 실마다 $32$개였다. 사람에게 이식한 장치는 회사 발표에 따르면 실 $64$가닥에 약 1000채널이다. 유연한 실은 조직을 덜 손상시키고, 두개골 안에서 뇌가 늘 조금씩 움직이는 것도 더 잘 견딘다. 실은 로봇이 혈관을 피해 꽂는다. 앞서 말했듯이 첫 논문에서는 원시 흐름~\eqref{eq:probedata}이 케이블로 나갔다. 회사 발표에 따르면 사람용 장치는 다르다. 본체는 완전히 피부로 덮이고, 스파이크는 안에서 골라내며, 데이터는 무선으로 나가고, 전력은 선 없이 들어온다.

첫 장치의 과제는 읽기다. 실을 팔 움직임을 계획하는 피질 부위에 꽂으면, 디코더가 스파이크를 화면의 커서 움직임으로 바꾼다. 팔이 마비된 사람이 움직임을 상상해서 컴퓨터를 조작한다. 발상 자체는 새롭지 않다. 전극 100개 남짓의 딱딱한 배열을 쓴 인터페이스는 오래전부터 커서를 조작하게 해 주었고~\cite{Hochberg2006}, 손으로 글씨 쓰는 움직임을 상상해 글을 쓰게 했으며~\cite{Willett2021}, 말하려는 시도를 분당 약 60단어의 속도로 화면의 글로 바꾸기까지 했다~\cite{Willett2023}. Neuralink의 새로움은 공학에 있다. 채널 수, 무선, 밀폐된 본체, 로봇 삽입, 즉 실험실 밖에서 몇 년 동안 지닐 수 있는 프로브를 만들려는 시도다. 우리가 아는 한 사람 대상 시험 결과는 주로 동료 심사 논문이 아니라 회사 자체가 발표했다. 회사는 특히 삽입 뒤 실 일부가 밀려나 작동하는 채널 수가 줄었다고 밝혔다. 디버깅에서는 흔한 골칫거리다. 보드에 대해 움직이는 프로브는 이미 다른 전선을 듣는다.

회사의 두 번째 방향은 쓰기다. 시력을 잃은 사람의 시각 피질에 신호를 넣는 시각 장치다. 이에 대해서는 아래 쓰기 절에서 다룬다.

\section{읽기: 수신자로서의 디코더}

인터페이스의 디코더는 명제~\ref{prop:reader}의 의미에서 수신자다. 메시지의 어느 부분을 읽을지 스스로 정한다. 사실상 모든 인터페이스는 \emph{무리의 발화율}을 읽는다. 각 채널의 스파이크를 수십 밀리초 창마다 세고, 의도한 움직임이 나오도록 고른 가중치로 더한다. 이것은 \ref{sec:rate}절의 무리 발화율 부호이며, 우연히 고른 것이 아니다. 창마다 스파이크가 몇 개뿐이면 뉴런 하나의 발화율은 정해지지 않지만, 100개를 합치면 이미 쓸 만하다.

정보를 얼마나 꺼낼 수 있을까? “생각 속의 손”으로 글쓰기는 분당 약 90자, 즉 초당 1.5자의 속도였다~\cite{Willett2021}. 글자당 5비트로 잡아도 초당 8비트가 안 된다. Neuralink 발표에 따르면 커서 조작은 초당 약 10비트다. 한편 상한~\eqref{eq:spikeentropy}은 뉴런 \emph{하나}에 초당 수십 비트, 100개면 수천 비트다. 인터페이스는 엿들은 뉴런이 전달할 수 있는 양의 작은 일부만 꺼낸다. 병목은 전선이 아니다. 무리가 독립 변수를 몇 개나 전달하는지(명제~\ref{prop:pooling}), 그리고 \ref{sec:code}장에서 보았듯이 아직 완전히 해독되지 않은 부호를 디코더가 얼마나 잘 이해하는지가 병목이다.

\ref{sec:motorlearning}장에서 익숙한 두 번째 특징도 있다. 디코더와 뇌는 서로에게서 배운다. 뇌는 커서가 어디로 갔는지 보고 오차를 받아 명령을 조정한다. 오차 전선을 통한 운동 학습과 같다. 엔지니어는 때때로 디코더의 가중치를 다시 맞춘다. 실 아래의 뉴런이 바뀌고 네트워크의 연결이 다시 학습되기 때문이다. 결국 서로에게 맞춰 가는 두 학습자가 생긴다. 이 쌍이 폭주하지 않게 하려면 학습 속도를 떼어 놓는 것이 편하다. 한쪽은 빠르게, 다른 쪽은 느리게 배우게 한다.

\begin{keyblock}
\begin{proposition}[뉴런 1000개짜리 프로브가 작동하는 이유]
\label{prop:probe}
보잘것없는 비율의 뉴런을 엿듣는 인터페이스가 움직임을 제어할 수 있는 것은, 무리의 활동이 저차원이고 그 잡음이 상관되어 있기 때문이다. 뉴런 수백 개가 공통 변수를 거의 다 전달한다. 같은 이유로 인터페이스는 초당 몇 비트밖에 꺼내지 못한다. 스파이크 흐름이 담을 수 있는 양이 아니라 과제의 독립 변수 개수만큼이다. 인터페이스의 성능은 측정되었다. 부호를 이해하면 디코더가 얼마나 좋아질지는 열린 문제다.
\end{proposition}
\end{keyblock}

\section{쓰기: 읽기보다 어렵다}

기계에 신호를 넣는 것은 읽는 것보다 어렵다. 전류를 흘리는 전극은 뉴런 하나가 아니라 주변의 모든 뉴런과 전선을 유형과 부호를 가리지 않고 흥분시킨다. \emph{어느} 뉴런이 \emph{언제} 발화했는지가 중요한 부호(\ref{sec:code}장)를 이런 전극으로 쓸 수는 없다. \emph{어디에}, \emph{얼마나 세게}를 대략 정할 수 있을 뿐이다.

시각에는 이것으로 어느 정도 충분하다. 시각 피질의 전극에 전류를 흘리면 사람은 시야의 특정 위치에서 빛나는 점을 본다. 점을 동시에, 최대 16개까지 켜자 시력을 잃은 사람이 몇몇 글자를 알아보고 물체의 경계를 구별했다~\cite{Fernandez2021}. 이것은 위치 부호이며, 쓸 수 있다. 그러나 \ref{sec:sensors}장을 떠올려 보자. 눈은 뇌로 픽셀이 아니라 압축된 설명, 즉 가장자리, 변화, 밝아짐과 어두워짐을 따로따로의 전선으로 보낸다. 피질에 “빛의 점”을 쓰는 것은 입력에서 전혀 다른 부호를 기다리는 기계에 픽셀을 쓰는 것이다. 그래서 인공 시각은 아직 전극 수로 기대할 만한 수준보다 거칠다. 전극이 필요 없는 시험 신호도 있다. 착시는 눈을 통해 특이한 입력을 넣어 기계를 평소 모드에서 벗어나게 하고, 오류 하나하나가 각자의 메커니즘을 가리킨다(\ref{sec:illusions}절).

\section{프로브가 보지 못하는 것}

전극은 주소 버스, 즉 \nt{GLUT} 뉴런의 스파이크를 듣고, 작은 \nt{GABA} 뉴런은 덜 잘 듣는다. 그러나 \ref{sec:serocomputer}--\ref{sec:acet}장에서 보았듯이 기계 전체의 동작 모드는 브로드캐스트 레지스터, 즉 \nt{SERO}, \nt{DOPA}, \nt{NORA}, \nt{ACET}의 농도가 정한다. 전극은 이것을 듣지 못한다. 원천 뉴런이 극히 적고, 그 신호는 프로브 옆의 스파이크가 아니라 부피 속의 농도이기 때문이다. 데이터 버스는 보지만 제어 레지스터를 보지 못하는 디버거는 늘 놀라게 된다. 같은 입력이 다른 응답을 내는데, 어딘가에서 $g$, $a$, $\tau_\gamma$가 바뀌었기 때문이다. 농도는 조직 안의 화학 센서로 잴 수 있지만~\cite{Bunin1998}, 인터페이스에서는 아직 이런 센서를 쓰지 않는다. 기계의 느린 두 번째 출력인 혈액 속 호르몬(\ref{sec:chemoutput}장)도 프로브의 시야 밖에 있다. 호르몬은 전극이 아니라 혈액 시료로 잰다.

프로브는 가중치도 보지 못한다. 그런데 기억과 학습한 것의 거의 전부가 바로 가중치에 저장된다. 가중치는 응답이 어떻게 바뀌는지로 짐작할 수 있을 뿐이다. 또 프로브는 사람이 느끼는 방식의 통증도 보지 못한다. \ref{sec:pain}장에서 보았듯이 경보 신호의 세기는 기계 자신, 즉 척수의 관문과 위에서 오는 조절기가 정하므로, 센서로는 얼마나 아픈지 읽을 수 없다.

\section{요약: 디버깅과 열린 문제}

\begin{table}[t]
\centering
\footnotesize
\setlength{\tabcolsep}{4pt}
\renewcommand{\arraystretch}{1.15}
\begin{tabular}{>{\raggedright}p{2.2cm}>{\raggedright}p{3.6cm}>{\raggedright}p{3.4cm}>{\raggedright\arraybackslash}p{4.0cm}}
\toprule
디버깅 기법 & 인터페이스가 하는 일 & 책이 설명하는 것 & 알려진 것 \\
\midrule
프로브 & 뉴런 수백--수천 개를 듣는다 & 저차원성과 상관된 잡음(\ref{sec:code}장) & 측정됨 \\
프로브에서 압축 & 원시 신호 대신 이벤트를 보낸다 & 스파이크 흐름의 용량~\eqref{eq:spikeentropy} & 공학적으로 해결됨 \\
읽기 & 무리 발화율 $\to$ 움직임, 글, 말 & 디코더는 수신자, 무리 발화율 부호 & 작동함; 초당 몇 비트 \\
공동 학습 & 뇌와 디코더가 서로에게 맞춘다 & 오차 전선(\ref{sec:motorlearning}장) & 관찰됨; 맞춤의 규칙은 연구 대상 \\
쓰기 & 전류가 시야에 점을 켠다 & 위치 부호는 쓸 수 있고 타이밍 부호는 못 쓴다(\ref{sec:sensors}장) & 점은 측정됨; 쓸모 있는 시각은 아직 없음 \\
레지스터 읽기 & 거의 하지 않는다 & 브로드캐스트 채널(\ref{sec:serocomputer}--\ref{sec:acet}장) & 화학 센서는 실험에 있음 \\
\bottomrule
\end{tabular}
\caption{기계 디버깅: 뇌-컴퓨터 인터페이스가 할 수 있는 일과 이 책의 각 장이 그에 대해 말하는 것.}
\label{tab:debugging}
\end{table}

표~\ref{tab:debugging}이 이 장을 요약한다. 뇌-컴퓨터 인터페이스는 이미 주소 버스에 프로브를 대고 초당 몇 비트를 읽을 수 있다. 이것이 애초에 작동한다는 사실은 \ref{sec:code}장의 저차원성과 상관된 잡음으로 설명된다. 기계에 쓰는 일은 거칠게만 할 수 있다. 기계의 입력 부호가 압축되어 있고 개별 전선에 주소가 매겨져 있기 때문이다. 그리고 기계를 학습 가능하고 조정 가능하게 만드는 제어 레지스터와 가중치는 아직 프로브에 거의 보이지 않는다.

\ref{sec:synthesis}장의 열린 문제 하나하나는 디버거의 과제이기도 하다. 어떤 부호를 읽을지는 프로브가 더 조밀해지고 디코더가 발화율만이 아니라 스파이크 타이밍을 쓸 줄 알게 되면 정해질 것이다. 다층 네트워크가 어떻게 학습하는지는 여러 층에 동시에 프로브를 대고 학습하면서 응답이 어떻게 바뀌는지 보면 확인할 수 있다. 브로드캐스트 채널이 무엇을 전달하는지는 전기 프로브에 화학 프로브가 더해지면 보일 것이다. 이 책은 기계의 동작과 모든 엔지니어가 아는 법칙으로 기계의 회로도를 읽으려는 시도다. 지금 만들어지고 있는 디버거는 이 회로도를 프로브로 확인할 수 있게 해 줄 것이다.

\section{연습문제}

\begin{exercise}[\lvlA]
\label{ex:17:1}
\emph{무선 채널 예산.} 피부를 통한 전송은 비트당 $1$~nJ이 들고, 송신기는 최대 $10$~mW를 쓸 수 있다. $N=1000$ 전극에서 원시 흐름~\eqref{eq:probedata}과 스파이크 흐름에 필요한 전력은 얼마인가? 원시 흐름이라면 비트당 에너지가 얼마여야 하는가?
\end{exercise}

\begin{exercise}[\lvlA]
\label{ex:17:2}
\emph{뉴런 100개에 든 비트.} 디코더는 $r=20$~spikes/s인 뉴런 $N$개의 스파이크를 $50\ms$ 창마다 더한다. 잡음의 쌍별 상관은 $c=0.1$이고, 신호는 무리 발화율을 (제곱평균으로) $30\%$ 바꾼다. $N=10$, $100$, $1000$, $N\to\infty$에서 창당 속도 $\tfrac12\log_2(1+\text{SNR})$을 추정하라. $c=0$, $N=1000$이면?
\end{exercise}

\begin{exercise}[\lvlB]
\label{ex:17:3}
\emph{키보드 인터페이스의 속도.} 인터페이스가 초당 1.5번 $N=32$개 글자 가운데 하나를 고른다. 의도한 글자는 확률 $P$로 나오고 오류는 똑같이 확률이 나뉜다. $P=0.99$, $0.94$, $0.8$에서 초당 몇 비트를 전달하는가? $P=0.94$에서 $10$~bit/s를 내려면 초당 몇 글자가 필요한가?
\end{exercise}

\begin{exercise}[\lvlB]
\label{ex:17:4}
\emph{모델링: 무리 디코더.} 뉴런 $N$개, $\lambda_i=[b+m\,\mathbf v\cdot\mathbf p_i]_+$, $b=20$, $m=15$~spikes/s, 창 $50\ms$, $\mathbf v$는 성분마다 표준편차 $0.5$. 모든 뉴런이 공통 잡음이 섞인 $\mathbf v+\boldsymbol\xi$, $\sigma_\xi=0.2$를 본다. 비편향 무리 벡터 디코더를 시뮬레이션하라. $N$이 얼마일 때 두 잡음이 같아지며, $N\to\infty$에서 창당 성분마다 몇 비트인가?
\end{exercise}

\begin{exercise}[\lvlB]
\label{ex:17:5}
\emph{두 학습자.} 뇌는 명령 $m=b\,v^*$를, 디코더는 $v=d\,m$을 낸다. $\langle v^{*2}\rangle=1$. 둘 다 시도마다 $\tfrac12(v-v^*)^2$에 대해 속도 $\eta$로 경사 하강을 한 걸음 한다. $b=1.2$, $d=1$에서 시작해 $\eta=0.8$, $1.1$, $1.5$, $2$로 시뮬레이션하라. 각각은 $\eta<2$에서 안정하다. 쌍은 $\eta$가 얼마일 때 안정한가?
\end{exercise}

\begin{exercise}[\lvlB]
\label{ex:17:6}
\emph{설계: 프로브 파이프라인.} 채널 $1024$개, 채널당 초당 $2$만 샘플, 뉴런 $10$~spikes/s, 무선 $1$~nJ/bit. 샘플의 백색 잡음에 대해 어떤 문턱값이면 채널당 거짓 스파이크가 $0.1$~Hz 이하인가? $20\ms$마다 스파이크 수를 $2$비트로 보낸다면 전력은 얼마이고, 원시 $10$비트 샘플보다 몇 배 작은가?
\end{exercise}

\begin{exercise}[\lvlC]
\label{ex:17:7}
\emph{인터페이스 속도의 한계.} 대역폭 $W=5$~Hz인 변수 $D=10$개에 대해 $\text{SNR}\approx0.9$(신호를 닮은 잡음)와 $90$(뉴런 $1000$개의 독립 잡음)일 때 $R\le DW\log_2(1+\text{SNR})$을 추정하라. 측정값은 약 $10$~bit/s다. 격차의 두 설명, 즉 신호를 닮은 잡음과 부호를 모른다는 것을 어떤 실험으로 구별하겠는가?
\end{exercise}
